\section{Background and Related Work}
\label{sec:related}

\subsection{Rubrics for OER quality}
The most direct prior work is the review of 14 OER quality rubrics by Yuan and Recker \cite{yuan2015}. It examined them along content, development process and application context, and found wide diversity: the number of indicators ranges from 3 to 42, five rubrics specify no rating scale, only three give a detailed scoring guide, only six report empirical testing and two target automated evaluation. The Learning Object Review Instrument (LORI) stands out in that review for combining a rating scale, a scoring guide, reported testing and several revisions. Other work asks users what they want. Clements and Pawlowski surveyed teachers who use OER repositories about how they re-use resources and perceive quality, and report that trust in resources, organizations and technologies matters in particular \cite{clements2012}. Bliss and colleagues studied college teacher and student perceptions of OER \cite{bliss2013}. A recent systematic review still lists quality assurance as a challenge \cite{adil2022}, and a meta-synthesis of the literature on open education practices shows how broadly the field is defined \cite{shareefa2023}. \emph{This paper differs} by taking the review's findings as requirements: a score must have a scale, anchors and evidence, and it must be tested before it is trusted.

\subsection{Reusability and metadata}
Reuse is a separate construct from teaching quality, yet rubrics mostly fold it in as one indicator or one emphasis. Work on learning objects proposed automatic metrics for the quality of metadata in digital repositories \cite{ochoa2009} and relevance-ranking metrics for learning objects \cite{ochoa2008}. In scientific data, the FAIR principles make findability, accessibility, interoperability and reusability explicit, machine-checkable properties \cite{wilkinson2016}. Openness itself has a defined set of permissions \cite{wiley5r,wiley2018}, and accessibility has a published standard \cite{wcag21}. \emph{This paper differs} by treating reusability as a co-equal half of the score, with its own five metrics, so that a resource can be shown to be strong in one half and weak in the other.

\subsection{Validity, reliability and rubric use}
A score is only as good as the argument that supports its use. Messick and Kane describe validation as building and testing an interpretation-and-use argument \cite{messick1995,kane2013}. Reviews of rubric research find that analytic rubrics with explicit descriptors can support consistent scoring, but evidence about validity is thinner than evidence about reliability \cite{jonsson2007,reddy2010}. Agreement between raters is reported with Cohen's kappa \cite{cohen1960}, intraclass correlation \cite{shrout1979,koo2016} and chance-corrected coefficients that behave better when agreement is high \cite{gwet2008}, interpreted with conventional benchmarks \cite{landis1977}. Measurement also changes behavior: once a metric becomes a target it stops being a good measure \cite{campbell1979,strathern1997}. \emph{This paper differs} by naming these methods as the required validation program for an OER score, and by flagging gaming as a design risk.

\subsection{AI agents as evaluators}
Large language models are used to grade and rate text, and agreement with humans can be high enough to be useful but is affected by position and verbosity biases \cite{zheng2023,wang2024}. Prompted judges that follow criteria align better with humans on some tasks \cite{liu2023}, and a survey catalogs strategies for making judges reliable \cite{gu2024}. Models also produce confident falsehoods, including fabricated citations \cite{ji2023}, and carry risks when used in education \cite{bender2021,kasneci2023,zawacki2019}. In education, one study compared crowdsourced and language-model-based applications of two standardized rubrics to 30 assessment items across six domains \cite{moore2024}, and the open-education community has begun to examine how generative AI relates to OER and open educational practices \cite{bozkurt2023}. \emph{This paper differs} by treating the agent as one possible reviewer inside a score whose validity is established separately, rather than as the source of the score's authority.
